\SourceLeaf{372}
\par 
Ex definitione numeri medii arithmetico-geometrici tam manifestum est, ut
explicatione ampliore iam non opus sit, sequens

\textsc{Theorema.}\quad \textit{Numerus medius inter terminos quoscunque correspondentes progressionum I, II idem est atque inter $a$ et $b$.}

\AGMArticle{3}

Quo clarius perspiciatur, quanta rapiditate series (I) et (II) ad dextram
versus limitem suum approximent, et quomodo versus laevam illa crescat, haec
decrescat, exempla quaedam hic sistimus:

\Needspace{20\baselineskip}\begin{center}\textit{Exemplum} 1.\quad $a=1,\ b=0{,}2$\end{center}
\begingroup\small\setlength{\arraycolsep}{0pt}
\[\begin{array}{r@{\,=\,}l@{\quad}r@{\,=\,}l}
{}^{\prime\prime\prime\prime}a&15{,}83795\,47919\,02&{}^{\prime\prime\prime\prime}b&0{,}00000\,00000\,00000\,00005\,7\\[6pt]
{}^{\prime\prime\prime}a&7{,}91897\,73959\,512&{}^{\prime\prime\prime}b&0{,}00000\,00013\,481\\[6pt]
{}^{\prime\prime}a&3{,}95948\,86986\,4971&{}^{\prime\prime}b&0{,}00010\,30955\,7682\\[6pt]
{}^{\prime}a&1{,}97979\,58971\,13271\,23927\,9&{}^{\prime}b&0{,}02020\,41028\,86728\,76072\,1\\[6pt]
a&1{,}00000\,00000\,00000\,00000\,0&b&0{,}20000\,00000\,00000\,00000\,0\\[6pt]
a'&0{,}60000\,00000\,00000\,00000\,0&b'&0{,}44721\,35954\,99957\,93928\,2\\[6pt]
a''&0{,}52360\,67977\,49978\,96964\,1&b''&0{,}51800\,40128\,22268\,36005\,0\\[6pt]
a'''&0{,}52080\,54052\,86123\,66484\,5&b'''&0{,}52079\,78709\,39876\,24344\,0\\[6pt]
a''''&0{,}52080\,16381\,12999\,95414\,3&b''''&0{,}52080\,16380\,99375\\[6pt]
a^{\mathrm V}&0{,}52080\,16381\,06187&b^{\mathrm V}&0{,}52080\,16381\,06187
\end{array}
\]\endgroup

Hic $a^{\mathrm V}$, $b^{\mathrm V}$ in $23^{\mathrm a}$ demum figura discrepant;
${}^{\mathrm V}b$ est minor quam $(\frac{1}{10})^{40}$;
${}^{\prime\prime\prime\prime}a$, ${}^{\mathrm V}a$, ${}^{\mathrm{VI}}a$ etc.
sensibiliter formant progressionem geometricam, cuius exponens $=2$.

\Needspace{20\baselineskip}\begin{center}\textit{Exemplum} 2.\quad $a=1,\ b=0{,}6$.\end{center}
\begingroup\small\setlength{\arraycolsep}{0pt}
\[\begin{array}{r@{\,=\,}l@{\quad}r@{\,=\,}l}
{}^{\prime\prime\prime\prime}a&14{,}35538\,2913&{}^{\prime\prime\prime\prime}b&0{,}00000\,00000\\[6pt]
{}^{\prime\prime\prime}a&7{,}17769\,14569\,307&{}^{\prime\prime\prime}b&0{,}00001\,73070\\[6pt]
{}^{\prime\prime}a&3{,}58885\,43819\,99831\,75712\,7&{}^{\prime\prime}b&0{,}01114\,56180\,00168\,24287\,3\\[6pt]
{}^{\prime}a&1{,}80000\,00000\,00000\,00000\,0&{}^{\prime}b&0{,}20000\,00000\,00000\,00000\,0\\[6pt]
a&1{,}00000\,00000\,00000\,00000\,0&b&0{,}60000\,00000\,00000\,00000\,0\\[6pt]
a'&0{,}80000\,00000\,00000\,00000\,0&b'&0{,}77459\,66692\,41483\,37703\,6\\[6pt]
a''&0{,}78729\,83346\,20741\,68851\,8&b''&0{,}78719\,58685\,06172\,16741\,6\\[6pt]
a'''&0{,}78724\,71015\,63456\,92796\,7&b'''&0{,}78724\,70999\\[6pt]
a''''&0{,}78724\,71007\,8&b''''&0.78724\,71007\,8
\end{array}
\]\endgroup
\SourceLeaf{373}
\Needspace{20\baselineskip}\begin{center}\textit{Exemplum} 3.\quad $a=1,\ b=0{,}8$.\end{center}
\begingroup\small\setlength{\arraycolsep}{0pt}
\[\begin{array}{r@{\,=\,}l@{\quad}r@{\,=\,}l}
{}^{\mathrm V}a&25{,}19190\,722&{}^{\mathrm V}b&0{,}00000\,00000\,0\\[6pt]
{}^{\prime\prime\prime\prime}a&12{,}59595\,36116\,78&{}^{\prime\prime\prime\prime}b&0{,}00000\,00133\,367\\[6pt]
{}^{\prime\prime\prime}a&6{,}29797\,68125\,07655\,42373\,4&{}^{\prime\prime\prime}b&0{,}00040\,98644\,58278\,08440\,9\\[6pt]
{}^{\prime\prime}a&3{,}14919\,33384\,82966\,75407\,2&{}^{\prime\prime}b&0{,}05080\,66615\,17033\,24592\,8\\[6pt]
{}^{\prime}a&1{,}60000\,00000\,00000\,00000\,0&{}^{\prime}b&0{,}40000\,00000\,00000\,00000\,0\\[6pt]
a&1{,}00000\,00000\,00000\,00000\,0&b&0{,}80000\,00000\,00000\,00000\,0\\[6pt]
a'&0{,}90000\,00000\,00000\,00000\,0&b'&0{,}89442\,71909\,99915\,87856\,4\\[6pt]
a''&0{,}89721\,35954\,99957\,93928\,2&b''&0{,}89720\,92687\,32734\\[6pt]
a'''&0{,}89721\,14321\,16346&b'''&0{,}89721\,14321\,13738\\[6pt]
a''''&0{,}89721\,14321\,15042&b''''&0{,}89721\,14321\,15042
\end{array}
\]\endgroup

\Needspace{20\baselineskip}\begin{center}\textit{Exemplum} 4.\quad $a=\surd 2,\ b=1$.\end{center}
\begingroup\small\setlength{\arraycolsep}{0pt}
\[\begin{array}{r@{\,=\,}l@{\quad}r@{\,=\,}l}
{}^{\prime\prime\prime\prime}a&19{,}17024\,37557\,69475\,31905\,0&{}^{\prime\prime\prime\prime}b&0{,}00000\,00009\,32560\,02627\,6\\[6pt]
{}^{\prime\prime\prime}a&9{,}58512\,18783\,51017\,67266\,3&{}^{\prime\prime\prime}b&0{,}00013\,37064\,06056\,69181\,0\\[6pt]
{}^{\prime\prime}a&4{,}79262\,77923\,78537\,18223\,7&{}^{\prime\prime}b&0{,}03579\,93323\,67652\,95745\,7\\[6pt]
{}^{\prime}a&2{,}41421\,35623\,73095\,04880\,2&{}^{\prime}b&0{,}41421\,35623\,73095\,04880\,2\\[6pt]
a&1{,}41421\,35623\,73095\,04880\,2&b&1{,}00000\,00000\,00000\,00000\,0\\[6pt]
a'&1{,}20710\,67811\,86547\,52440\,1&b'&1{,}18920\,71150\,02721\,06671\,7\\[6pt]
a''&1{,}19815\,69480\,94634\,29555\,9&b''&1{,}19812\,35214\,93120\,12260\,7\\[6pt]
a'''&1{,}19814\,02347\,93877\,20908\,3&b'''&1{,}19814\,02346\,77307\,20579\,8\\[6pt]
a''''&1{,}19814\,02347\,35592\,20744\,1&b''''&1{,}19814\,02347\,35592\,20743\,9
\end{array}
\]\endgroup

\AGMArticle{4}

Habeantur praeter progressiones
\[
\begin{array}{l}
.\;.\;.{}^{\prime\prime\prime}a,\;{}^{\prime\prime}a,\;{}^{\prime}a,\; a,\; a',\; a'',\; a'''.\;.\;.\\[6pt]
.\;.\;.{}^{\prime\prime\prime}b,\;{}^{\prime\prime}b,\;{}^{\prime}b,\; b,\; b',\; b'',\; b'''.\;.\;.
\end{array}
\]
duae aliae
\[
\begin{array}{l}
.\;.\;.{}^{\prime\prime\prime}c,\;{}^{\prime\prime}c,\;{}^{\prime}c,\; c,\; c',\; c'',\; c'''.\;.\;.\\[6pt]
.\;.\;.{}^{\prime\prime\prime}d,\;{}^{\prime\prime}d,\;{}^{\prime}d,\; d,\; d',\; d'',\; d'''.\;.\;.
\end{array}
\]
simili modo formatae; supponamusque, duos terminos in posterioribus duobus in
prioribus \textit{proportionales} esse, e.\ g. $a:b=c:d$, sive $a:c=b:d=1:n$.
Tunc omnes termini in I ad terminos in III, omnesque in II ad terminos in IV
(similiter relative ad $a$, $b$, $c$, $d$ siti ac sitos) in eadem ratione erunt, puta
\SourceLeaf{374}
\[
\begin{gathered}c'=na',\quad c''=na'',\quad c'''=na''',\\ {}^{\prime}c={}^{\prime}an,\quad d^{\mathrm V}=nb^{\mathrm V},\quad{}^{\mathrm V}d={}^{\mathrm V}bn\end{gathered}
\]
Hinc facile deducitur, etiam limitem serierum I, II, fore ad limitem serierum
III, IV ut $1$ ad $n$, sive generaliter $\mathrm M(na,nb)=n\mathrm M(a,b)$.
Erit itaque generaliter $\mathrm M(a,b)=a\mathrm M(1,\frac{b}{a})=b\mathrm M(\frac{a}{b},1)$.

\AGMArticle{5}

\textsc{Problema.}\quad \textit{Exprimere medium arithmetico-geometricum inter numerum unitate maiorem $1+x$ et unitatem, per seriem secundum potestates ipsius $x$ progredientem.}

\textit{Sol.}\quad Quum $\mathrm M(1,1)=1$, supponamus
\[
\mathrm M(1+x,1)=1+h'x+h''x^2+h'''x^3+h''''x^4+\text{etc.}
\]
ita ut $h'$, $h''$, $h'''$, $h''''$ sint coëfficientes constantes ab $x$ non pendentes.
Sit $x=2t+tt$, eritque
\[
\begin{aligned}&\mathrm M(1+x,1)=\mathrm M(1+t+\tfrac12tt,\ 1+t)\\&=(1+t)\mathrm M(1+\frac{\tfrac12tt}{1+t},1).\end{aligned}
\]
Quare habebitur
\[
\begin{aligned}
&1+h'(2t+tt)+h''(2t+tt)^2+h'''(2t+tt)^3+\text{etc.}\\[6pt]
={}&1+t+h'(\tfrac12tt)+h''\frac{\tfrac14t^4}{1+t}+h'''\frac{\tfrac18t^6}{(1+t)^2}+\text{etc.}
\end{aligned}
\]
Hinc prodeunt aequationes
\[
\begin{aligned}
2h'&=1\\[6pt]
4h''+h'&=\tfrac12h'\\[6pt]
8h'''+4h''&=0\\[6pt]
16h''''+12h'''+h''&=\tfrac14h''\\[6pt]
32h^{\mathrm V}+32h''''+6h'''&=-\tfrac14h''\\[6pt]
64h^{\mathrm{VI}}+80h^{\mathrm V}+24h''''+h'''&=\tfrac14h''+\tfrac18h'''\\[6pt]
128h^{\mathrm{VII}}+192h^{\mathrm{VI}}+80h^{\mathrm V}+8h''''&=-\tfrac14h''-\tfrac28h'''\\[6pt]
256h^{\mathrm{VIII}}+448h^{\mathrm{VII}}+240h^{\mathrm{VI}}+40h^{\mathrm V}+h''''&=\tfrac14h''+\tfrac38h'''+\tfrac1{16}h''''
\end{aligned}
\]
etc.\quad unde fit
\[
\begin{gathered}h'=\tfrac12,\quad h''=-\tfrac1{16},\quad h'''=\tfrac1{32},\\ h''''=-\tfrac{21}{1024},\quad h^{\mathrm V}=\tfrac{31}{2048},\quad h^{\mathrm{VI}}=-\tfrac{195}{16384}\end{gathered}
\]
Quare
\[
\begin{aligned}&\mathrm M(1+x,1)=1+\tfrac12x-\tfrac1{16}xx+\tfrac1{32}x^3\\&\quad-\tfrac{21}{1024}x^4+\tfrac{31}{2048}x^5-\tfrac{195}{16384}x^6\ \text{etc.}\end{aligned}
\]
Ceterum nullo negotio perspicitur, medium inter $1$ et numerum unitate minorem
\SourceLeaf{375}
\(1-x\) fore \(1-h'x+h''xx-h'''x^3+\text{etc.}=1-\frac12x-\frac1{16}xx-\frac1{32}x^3-\frac{21}{1024}x^4-\text{etc.}\)
Quum hi co\"efficientes legem obviam non exhibeant, has series praetergredimur,
aliamque viam tentamus, quae successum feliciorem praestabit.

\AGMArticle{6}

\textsc{Problema.} \textit{Exprimere medium ar. g. inter \(1+x\) et \(1-x\) per seriem secundum potestates ipsius \(x\) progredientem.}

\textit{Sol.} Quum habeatur
\[
\mathrm{M}(1+x,1-x)=(1-x)\mathrm{M}(1+\frac{2x}{1-x},1)
\]
statim habetur e serie art. praec. substituendo ibi \(\frac{2x}{1-x}\) pro \(x\)
\[
\begin{aligned}&\mathrm{M}(1+\frac{2x}{1-x},1)=1+x+\frac34xx+\frac34x^3\\&\quad+\frac{43}{64}x^4+\frac{43}{64}x^5+\frac{161}{256}x^6+\text{etc.}\end{aligned}
\]
atque hinc
\[
\mathrm{M}(1+x,1-x)=1-\frac14xx-\frac5{64}x^4-\frac{11}{256}x^6+\text{etc.}
\]

Co\"efficientes huius seriei etiam independenter a serie art. praec. per methodum sequentem erui possunt. Ponatur \(x=\frac{2t}{1+tt}\) eritque
\[
\begin{aligned}&\mathrm{M}(1+x,1-x)=\mathrm{M}(\frac{1-tt}{1+tt},1)\\&=\frac1{1+tt}\mathrm{M}(1+tt,1-tt)\end{aligned}
\]
Quare statuendo
\[
\mathrm{M}(1+x,1-x)=1+\alpha xx+\beta x^4+\gamma x^6+\delta x^8+\text{etc.}
\]
(nam potestates ipsius \(x\) cum exponente impari non adesse sponte patet) habebitur
\[
\begin{aligned}
&(1+tt)\{1+\alpha(\frac{2t}{1+tt})^2+\beta(\frac{2t}{1+tt})^4\\&\quad+\gamma(\frac{2t}{1+tt})^6+\delta(\frac{2t}{1+tt})^8+\text{etc.}\}\\[6pt]
&\qquad=1+\alpha t^4+\beta t^8+\gamma t^{12}+\delta t^{16}+\text{etc.}
\end{aligned}
\]
Hinc prodeunt aequationes
\[
\begin{aligned}
1+4\alpha&=0\\&\qquad\text{unde }\alpha=-\frac14\\[6pt]
-4\alpha+16\beta&=\alpha\\&\qquad\beta=-\frac5{64}\\[6pt]
4\alpha-48\beta+64\gamma&=0\\&\qquad\gamma=-\frac{11}{256}\\[6pt]
-4\alpha+96\beta-320\gamma+256\delta&=\beta\\&\qquad\delta=-\frac{29}{1024}-\frac5{16384}=-\frac{469}{16384}\\[6pt]
&\text{etc.}\qquad\text{etc.}
\end{aligned}
\]
In hac quoque serie co\"efficientes legi simplici non subiecti sunt: at si unitas per
illam seriem dividitur, prodit
\SourceLeaf{376}
\[
\begin{aligned}&\frac1{\mathrm{M}(1+x,1-x)}=1+\frac14xx+\frac9{64}x^4+\frac{25}{256}x^6\\&\quad+\frac{1225}{16384}x^8+\text{etc.}\end{aligned}
\]
ubi primo aspectu videmus, co\"efficientes esse quadrata radicum
\(\frac12,\allowbreak \ \frac12.\frac34,\allowbreak \ \frac12.\frac34.\frac56,\allowbreak \ \frac12.\frac34.\frac56.\frac78\),
adeoque secundum legem persimplicem progredi. Sed methodus, per
quam ad hanc conclusionem pulcherrimam pervenimus, inductionis tantummodo
vim habet; quam ad certitudinis gradum evehere in disquisitionibus sqq. nobis
proponimus.

\AGMArticle{7}

Supponendo
\[
\frac1{\mathrm{M}(1+x,1-x)}=1+Axx+Bx^4+Cx^6+\text{etc.}
\]
atque ut in art. praec. \(x=\frac{2t}{1+tt}\) habemus
\[
\begin{aligned}
&1+A(\frac{2t}{1+tt})^2+B(\frac{2t}{1+tt})^4\\&\quad+C(\frac{2t}{1+tt})^6+D(\frac{2t}{1+tt})^8+\text{etc.}\\[6pt]
&=1+tt+At^4+At^6+Bt^8+Bt^{10}\\&\quad+Ct^{12}+Ct^{14}+Dt^{16}+Dt^{18}+\text{etc.}
\end{aligned}
\]
unde emergunt aequationes
\[
\begin{aligned}
4A&=1\\[6pt]
-8A+16B&=A\\[6pt]
12A-64B+64C&=A\\[6pt]
-16A+160B-384C+256D&=B\qquad\text{etc.}
\end{aligned}
\]
atque hinc \(A=\frac14,\allowbreak \ B=\frac14.\frac9{16},\allowbreak \ C=\frac14.\frac9{16}.\frac{25}{36},\allowbreak \ D=\frac14.\frac9{16}.\frac{25}{36}.\frac{49}{64}\), ut supra:
sed legis ratio hinc operosius deduceretur: quare methodum sequentem praeferimus. Ex aequatione
\[
\begin{aligned}&\frac{2t}{1+tt}+A(\frac{2t}{1+tt})^3+B(\frac{2t}{1+tt})^5+\text{etc.}\\&=2t(1+At^4+Bt^8.\;.\;.)\end{aligned}
\]
demanant aequationes sequentes:
\[
\begin{array}{@{}rcl@{\quad}r@{}}
1&=&1&[1]\\[6pt]
0&=&1-4A&[2]\\[6pt]
A.&=&1-12A+16B&[3]\\[6pt]
0&=&1-24A+80B-64C&[4]\\[6pt]
B&=&1-40A+240B-448C+256D&[5]\\[6pt]
0&=&1-60A+560B-1792C+2304D-1024E&[6]
\end{array}
\]
\SourceLeaf{377}
ubi co\"efficientes facile subiiciuntur formulae generali: scilicet aequatio \(n^{\mathrm{ta}}\) erit
\[
\begin{aligned}
M={}&1-4A\times\frac{n.n-1}{1.2}\\&\quad+16B\times\frac{n+1.n.n-1.n-2}{1.2.3.4}\\&\quad-64C\times\frac{n+2.n+1.n.n-1.n-2.n-3}{1.2.3.4.5.6}\\[6pt]
&+256D\times\frac{n+3.n+2.n+1.n.n-1.n-2.n-3.n-4}{1.2.3.4.5.6.7.8}-\text{etc.}
\end{aligned}
\]
ubi \(M\) erit vel \(=0\) (quando \(n\) par), vel aequalis termino \(\frac12(n+1)^{\mathrm{to}}\) seriei
\(1,\ A,\ B,\ C,\ D\) etc. (quando \(n\) impar). Iam ex his aequationibus sequentes novas deducimus\HistoricalNote{27}.
\[
\begin{aligned}
\relax [2],\quad&0=1-4A\\[6pt]
4[3]-[1],\quad&4A-1=3-48A+64B\\[6pt]
9[4]-4[2],\quad&0=5-200A+720B-576C\\[6pt]
16[5]-9[3],\quad&16B-9A=7-532A+3696B-7168C\\&\quad+4096D\\[6pt]
25[6]-16[4],\quad&0=9-1116A+12720B\\&\quad-43776C+57600D-25600E
\end{aligned}
\]
etc., ubi co\"efficientes legi generali facile subiiciuntur. Scilicet aequatio \(n^{\mathrm{ta}}\) erit
\[
\begin{aligned}
&nnN-(n-1)^2L\\[6pt]&=(2n-1)(1-4A\times\frac{3nn-3n+2}{1.2}
\\&\quad+16B\times\frac{n.n-1.5nn-5n+6}{1.2.3.4}\\[6pt]
&-16C\times\frac{n+1.n.n-1.n-2.7nn-7n+12}{1.2.3.4.5.6}\\[6pt]
&+64D\times\frac{n+2.n+1.n.n-1.n-2.n-3.9nn-9n+20}{1.2.3.4.5.6.7.8}-\text{etc.})\quad\text{\textsuperscript{**)}}
\end{aligned}
\]\begingroup\renewcommand{\thefootnote}{**)}\footnotetext{\csname HistoricalNoteText28\endcsname}\endgroup
ubi factorum progressio obvia est (puta praeter factores simplices in singulos co\"efficientes ingreditur factor duplex talis \(knn-kn+\frac14(kk-1)\)). Hae aequationes
simplicius sequenti modo exhibentur, singularum membris ad dextram in binas
partes discerptis (praeter aequ. primam, quae immutata retinetur):
\[
\begin{aligned}
0&=1-4A\\[6pt]
4A-1&=\{3-36A\\[6pt]
&\phantom{=\{}-12A+64B\}\\[6pt]
0&=\{5-180A+400B\\[6pt]
&\phantom{=\{}-20A+320B-576C\}\\[6pt]
16B-9A&=\{7-504A+2800B-3136C\\[6pt]
&\phantom{=\{}-28A+896B-4032C+4096D\}\\[6pt]
0&=\{9-1080A+10800B-28224C+20736D\\[6pt]
&\phantom{=\{}-36A+1920B-15552C+36864D-25600E\}
\end{aligned}
\]
\SourceLeaf{378}
scilicet generaliter aequatio \(n^{\mathrm{ta}}\)
\[
\begin{aligned}
&nnN-(n-1)^2L\\[6pt]
&=(2n-1)\Biggl\{\begin{aligned}
1&-3.4A\frac{n. n-1}{1.2}
\\&\quad+5.16B\frac{n+1. n. n-1. n-2}{1.2.3.4}
\\&\quad-7.64C\frac{n+2. n+1.\;.\;.\;. n-3}{1.2.\;.\;.6}\\[6pt]
&-4A\frac{2}{1.2}
\\&\quad+16B\frac{n. n-1.16}{1.2.3.4}
\\&\quad-64C\frac{n+1. n. n-1. n-2.54}{1.2.3.4.5.6}
\end{aligned}\\[6pt]
&\qquad\begin{aligned}
&.\;.\;.\;.\;.\pm k.2^{k-1}K
\frac{n+\frac{k-3}{2}. n+\frac{k-5}{2}. n+\frac{k-7}{2}.\;.\;. n-\frac{k-1}{2}}
{1.2.3.\;.\;.\;. k-1}.\;.\;.\\[6pt]
&.\;.\;.\;.\;.\pm 2^{k-1}K
\frac{n+\frac{k-5}{2}. n+\frac{k-7}{2}.\;.\;.\;. n-\frac{k-3}{2}.\frac14(k-1)^3}
{1.2.\;.\;.\;.\;.\;.\;.\;.\;. k-1}.\;.\;.
\end{aligned}\Biggr\}
\end{aligned}
\]
designante \(k\) indefinite quemvis imparem, sint ita
\[
\begin{aligned}
0&=1-4A\\[6pt]
4A-1&=3(1-4A)-4(9A-16B)\\[6pt]
0&=5(1-4A)-20(9A-16B)+16(25B-36C)\\[6pt]
16B-9A&=7(1-4A)-56(9A-16B)+112(25B-36C)\\&\quad-64(49C-64D)\\[6pt]
0&=9(1-4A)-120(9A-16B)\\&\quad+432(25B-36C)-576(49C-64D)\\[6pt]
&\qquad+256(81D-100E)
\end{aligned}
\]
etc. et generaliter aequ. \(n^{\mathrm{ta}}\)
\[
\begin{aligned}
&nnN-(n-1)^2L\\[6pt]
&=(2n-1)(1-4A)\\&\quad-4\frac{2n-1. n. n-1}{1.2.3}(9A-16B)\\&\quad+16\frac{2n-1. n+1. n. n-1. n-2}{1.2.3.4.5}(25B-36C)\\[6pt]
&\qquad-64\frac{2n-1. n+2. n+1. n. n-1. n-2. n-3}{1.2.3.4.5.6.7}(49C-64D)+\text{etc.}
\end{aligned}
\]
ubi legis obviae universalitas e calculo sponte demanat. Hinc vero perspicuum
est, fieri necessario
\[
\begin{gathered}0=1-4A,\quad 0=9A-16B,\quad 0=25B-36C,\\[6pt]
0=49C-64D,\quad 0=81D-100E\end{gathered}
\]
\textit{etc. in inf.} adeoque
\[
\begin{gathered}A=\tfrac14,\\ B=\tfrac14.\tfrac9{16},\\[6pt]
C=\tfrac14.\tfrac9{16}.\tfrac{25}{36},\\[6pt]
D=\tfrac14.\tfrac9{16}.\tfrac{25}{36}.\tfrac{49}{64},\\[6pt]
E=\tfrac14.\tfrac9{16}.\tfrac{25}{36}.\tfrac{49}{64}.\tfrac{81}{100}\end{gathered}
\]
\textit{et sic porro in infinitum.} Q.~E.~D.
\SourceLeaf{379}
\AGMArticle{8}
Si statuimus
\[
1+\tfrac14xx+\tfrac14.\tfrac9{16}x^4
+\tfrac14.\tfrac9{16}.\tfrac{25}{36}x^6+\text{etc.}=y
\]
fit
\[
\begin{aligned}&\tfrac12xx+\tfrac14.\tfrac94x^4\\&\quad+\tfrac14.\tfrac9{16}.\tfrac{25}{6}x^6\\&\quad+\tfrac14.\tfrac9{16}.\tfrac{25}{36}.\tfrac{49}{8}x^8
+\text{etc.}=\frac{x\mathrm{d}y}{\mathrm{d}x}\end{aligned}
\]
atque
\[
\begin{aligned}&xx+\tfrac14.9x^4+\tfrac14.\tfrac9{16}.25x^6\\&\quad+\tfrac14.\tfrac9{16}.\tfrac{25}{36}.49x^8+\text{etc.}\\&=\frac{xx\mathrm{dd}y}{\mathrm{d}x^2}+\frac{x\mathrm{d}y}{\mathrm{d}x}\end{aligned}
\]
unde sponte sequitur
\[
\frac{xx\mathrm{dd}y}{\mathrm{d}x^2}+3\frac{x\mathrm{d}y}{\mathrm{d}x}+y
=\frac1{xx}\left(xx\frac{\mathrm{dd}y}{\mathrm{d}x^2}+\frac{x\mathrm{d}y}{\mathrm{d}x}\right)
\]
sive
\[
(x^3-x)\frac{\mathrm{dd}y}{\mathrm{d}x^2}
+(3xx-1)\frac{\mathrm{d}y}{\mathrm{d}x}+xy=0
\]
Hoc itaque modo media nostra arithmetico-grometrica ad quantitates integrales
revocata sunt, solutionemque particularem huiusce aequationis differentio-differentialis subministrant.

Eiusdem aequationis int. compl. est
\(\dfrac{\mathfrak{A}}{\mathrm{M}(1+x,\,1-x)}+\dfrac{\mathfrak{B}}{\mathrm{M}(1,x)}\).

Sit \(\varphi\) angulus indefinitus, eritque valor integralis
\(\int\cos\varphi^2\mathrm{d}\varphi\), a \(\varphi=0\) usque ad
\(\varphi=\pi\), ut vulgo notum est, \(=\tfrac12\pi\); eodem modo fit
valor integralis \(\int\cos\varphi^4\mathrm{d}\varphi\) inter eosdem limites
\(=\tfrac12.\tfrac34\pi\); valor integralis
\(\int\cos\varphi^6\mathrm{d}\varphi=\tfrac12.\tfrac34.\tfrac56\pi\)
etc. --- denique, ut sponte patet, \(\int\mathrm{d}\varphi=\pi\).
Hinc perspicuum est, valorem integralis
\[
\begin{aligned}&\int\mathrm{d}\varphi\times
\bigl(1+\tfrac12x^2\cos\varphi^2\\&\quad+\tfrac12.\tfrac34x^4\cos\varphi^4\\&\quad+\tfrac12.\tfrac34.\tfrac56x^6\cos\varphi^6+\text{etc.}\bigr)\end{aligned}
\]
sive huius \(\int\dfrac{\mathrm{d}\varphi}{\surd(1-xx\cos\varphi^2)}\),
fieri \(=\pi y\), si sumatur a \(\varphi=0\) usque ad \(\varphi=\pi\),
spectando quantitatem \(x\) tamquam constantem.

Quodsi functio \(\dfrac1{\surd(1-xx\cos\varphi^2)}\) in seriem talem
evolvi supponatur,
\[
P+2Q\cos2\varphi+2R\cos4\varphi+2S\cos6\varphi+\text{etc.}
\]
ita ut co\"efficientes \(P,Q,R,S\) etc. a sola \(x\) pendeant: valor
integralis supra traditi completus erit
\[
P\varphi+Q\sin2\varphi+\tfrac12R\sin4\varphi+\tfrac13S\sin6\varphi
+\text{etc.}+\text{Const.}
\]
\SourceLeaf{380}
adeoque valor intra limites ante allatos, $=P\varphi$, unde $y=P$. Iam observamus,
quum sit $\frac{1}{y}=\mathrm{M}(1+x,1-x)$ fieri quoque
$\frac{1}{y}=\mathrm{M}(1,\surd(1-xx))$. Hinc facillime derivatur sequens
theorema generalius: Si expressio talis
$\frac{\alpha}{\surd(\beta-\gamma\cos\varphi^{2})}=W$
in seriem secundum cosinus angulorum $2\varphi$, $4\varphi$, $6\varphi$ etc.
progredientem evolvatur, cuius terminus constans designetur per $P$; valores maximus
et minimus ipsius $W$ autem (puta $\frac{\alpha}{\surd(\beta-\gamma)}$ et
$\frac{\alpha}{\surd\beta}$) denotentur per $v,v'$: erit $\frac{1}{P}$ medium arithmetico-geometricum inter $\frac{1}{v}$ et $\frac{1}{v'}$. Proxime quidem
ratiocinia praecedentia pro eo tantummodo casu valent, ubi $\gamma$ est quantitas
positiva: sed facillime ad eum quoque casum extenduntur, ubi $\gamma$ est negativus.
In hocce enim casu fit $W=\frac{\alpha}{\surd(\beta-\gamma+\gamma\cos\psi^{2})}$
ponendo $\psi=90-\varphi$; unde $\frac{1}{P}$ med. inter
$\frac{\surd\beta}{\alpha}$ et $\frac{\surd(\beta-\gamma)}{\alpha}$. Ut supra.
Ceterum nullo negotio patet, idem theorema sine ulla mutatione etiam ad expressiones tales $\frac{\alpha}{\surd(\beta+\gamma\cos\varphi)}$ extendi, puta ut
$\frac{1}{\text{Term. Const.}}$ fiat
\[
\begin{gathered}\\=\mathrm{M}\left(\frac{1}{\text{valor max.}},\frac{1}{\text{valor minim.}}\right)\\=\mathrm{M}\left(\frac{\surd(\beta-\gamma)}{\alpha},
                 \frac{\surd(\beta+\gamma)}{\alpha}\right)\end{gathered}
\]
huiusmodi enim functiones reducentur ad formam
$\frac{\alpha}{\surd(\beta-\gamma+2\gamma\cos\psi^{2})}$, ei de qua modo diximus
prorsus similem, si scribatur $\varphi=2\psi$\HistoricalNote{29}.

Denique monemus, in sequentibus demonstrationem multo generaliorem eorundem
theorematum ex principiis magis genuinis datum iri; praetereaque mox etiam omnes
reliquos coëfficientes $Q,R,S$ etc. per methodos aeque expeditas eruere docebimus.
Hoc loco haec disquisitio eam quoque utilitatem praestabit, ut iis, qui dum
veritatum aeternarum sublimitatem atque divinam venustatem non sapiunt, earum
pretium ex solo usu, qui inde in partes matheseos applicatae redundare potest
aestimare noverunt,---hasce investigationes cariores reddat. Quantae enim utilitatis
sit evolutio coëfficientium $P,Q,R,S$ etc. tam rapida, ut ea quae ex his principiis
promanat, in astronomia physica sive theoria perturbationum planetarum, nemo ignorat.


\par\begin{center}\rule{.15\linewidth}{.4pt}\end{center}
\clearpage\fancyhead[L]{\small DE FUNCTIONIBUS TRANSSCENDENTIBUS}
\pdfbookmark[0]{PARS II. DE FUNCTIONIBUS TRANSSCENDENTIBUS}{agm-pars2}
\SourceLeaf{381}
\begin{center}
\textbf{PARS II.}\par
\textbf{DE FUNCTIONIBUS TRANSSCENDENTIBUS QUAE EX DIFFERENTIATIONE}\par
\textbf{MEDIORUM ARITHMETICO-GEOMETRICORUM ORIUNTUR.}
\end{center}
\par\begin{center}\rule{.15\linewidth}{.4pt}\end{center}
\AGMArticle{9}
Sint
\[
\left\{\begin{array}{llllll}
x,&x',&x'',&x'''\; .\; .\; .\; .&&\text{(I)}\\[6pt]
y,&y',&y'',&y'''\; .\; .\; .\; .&&\text{(II)}
\end{array}\right.
\]
series perinde formatae ut series in art. 1., puta ut quivis terminus in
$\left\{\begin{array}{c}\mathrm{I}\\\mathrm{II}\end{array}\right\}$ sit medium
$\left\{\begin{array}{c}\text{arithmeticum}\\\text{geometricum}\end{array}\right\}$
inter duos terminos praecedentes, adeoque $x^{\infty}=y^{\infty}=\mathrm{M}(x,y)$:
patetque, omnes has quantitates et proin etiam ipsarum limitem esse functiones
duarum variabilium $x,y$, atque variari, simulacque harum alteruter aut uterque
mutationem patiatur. Hic nobis de solis mutationibus infinite parvis sermo erit.
Fit itaque
\[
\begin{aligned}&\mathrm{d}x'=\tfrac12\mathrm{d}x+\tfrac12\mathrm{d}y,\\&
\mathrm{d}y'=\tfrac12\surd\frac{y}{x}. \mathrm{d}x
+\tfrac12\surd\frac{x}{y}. \mathrm{d}y\\&=\tfrac12\frac{y'}{x}\mathrm{d}x+\tfrac12\frac{y'}{y}\mathrm{d}y\end{aligned}
\]
et. perinde
\[
\begin{aligned}&\mathrm{d}x''=\tfrac12\mathrm{d}x'+\tfrac12\mathrm{d}y',\\&
\mathrm{d}y''=\tfrac12\frac{y''}{x'}. \mathrm{d}x'
+\tfrac12\frac{y''}{y'}\mathrm{d}y'\quad\text{etc.}\end{aligned}
\]
Hoc modo differentialia omnium terminorum in I et II per differentialia ipsarum
$x,y$ adiumento substitutionum exhiberi possent: sed lex progressiones magnopere
obscura hinc prodiret. Quam ut clarissime ob oculos producamus, sequentibus
scribendi compendiis utemur: Scribemus
